\section{Introduction}
Endorsements are a powerful tool for political elites interested in shaping primary elections and influencing candidate selection. Public support for a candidate serves as a signal to voters, shaping preferences and governing candidate viability. However, endorsements are not merely a signal of information to voters. They also play an integral role in intraparty dynamics by positioning and aligning internal factions with candidates. Party elites and leaders face a dilemma between endorsing the candidate most closely aligned with them in ideology and endorsing the candidate most likely to win.

An endorsement can signal policy commitments and further situate an elite within their party faction. However, a risk is involved in endorsing a candidate who may lose, jeopardizing the elite’s future political power. Conversely, endorsing the winning candidate associates them with the dominant coalition and secures influence in the new administration—often making it worthwhile to compromise on ideological preferences. The extent of this trade-off depends on the degree of competition between candidates and within-party polarization. This paper's first contribution is examining how ideological motivations and political incentives interact to shape endorsement behavior.

However, elite decisions are not made in isolation. Rather than treating them as a monolithic entity, I argue that the political elite is a composition of individuals with overlapping but distinct objectives. Some endorsers prioritize ideological purity. Others are more concerned with securing influence, positioning themselves for appointments, or shaping broader party narratives. The resulting strategic environment depends on both individual behavior and how elites respond to one another.

Existing research often treats the parties as uniform actors coordinating their influence on the nomination process. One of the most influential works on the political elite \textit{The Party Decides: Presidential Nominations Before and After Reform} \parencite{Cohen2008} argues that after the political reforms of the 1970s, parties - acting as a collective - use endorsements as a mechanism for signaling its preferred candidate. This perspective of the elite aligns with broader theories of party control, which suggest parties are long-term coalitions with interests in maintaining internal cohesion (\cite{aldrich1995}; \cite{aldrich2011}; \cite{Bawn2012}). In this view, party elites are expected to consolidate support behind a frontrunner to enhance electability and avoid intra-party fragmentation. However, while party unity is often a stated objective, the strategic incentives of individual elites may not always align with this collective goal. Rather than acting as a uniform bloc, party elites are heterogeneous actors with competing personal and factional interests \parencite{Noel2014}. If endorsements are driven by individual calculations rather than a unified party strategy, an important question follows: \textit{Are endorsements mutually reinforcing, or do they dilute each other’s impact?}

The key contribution of this paper is to develop a formal model that captures the strategic interaction between multiple endorsers and examines how political incentives shape the equilibrium pattern of endorsements. The model allows for both coordination and divergence in endorsements, depending on how political rewards for endorsing the winning candidate are structured.  Specifically, when political rewards are constant regardless of the number of endorsers, coordination emerges whenever ideological alignment between elites is strong enough. In this setting, endorsements are strategic complements—an elite’s endorsement of a candidate reduces the threshold of ideological alignment for others to endorse the same candidate. However, when political rewards depend on the number of endorsers, elites have an incentive to support different candidates to maximize their influence over the nomination. In this case, endorsements may act as strategic substitutes where elites diverge in support to secure greater political returns.

The model captures these dynamics by formalizing an endorsement game in which party elites decide whom to endorse based on both ideological affinity and strategic considerations. Each candidate is characterized by an exogenous quality level that reflects their electability, which, combined with endorsements, determines their likelihood of securing the nomination. Endorsements affect voter perceptions but are subject to diminishing marginal returns. The model introduces uncertainty in the nomination process, where the probability of a candidate’s success is influenced by both the number of endorsements received and an exogenous shock. Elites must weigh the expressive benefit of endorsing an ideologically aligned candidate against the strategic advantage of backing the most electable candidate.

This paper adds to the literature on elite influence in primaries by extending previous models of endorsements in two important ways. While prior research has considered the role of ideology in endorsements, this paper formally incorporates the trade-off between ideological alignment and political incentives in a multi-endorser framework. In addition, by allowing the strategic incentives of endorsers to vary based on the political reward structure, this paper highlights how different institutional settings can shape the strategic nature of endorsements. Instead of treating endorsements as purely informative or always reinforcing, I endogenize the relationship between endorsements to clarify the conditions under which they act as complements versus substitutes. The findings provide theoretical intuition for why early endorsements often appear fragmented but later endorsements tend to coalesce as elites coordinate their support.

The remainder of the paper is structured as follows. Section II provides an overview of the related literature, focusing on both theoretical and empirical research related to political endorsements. Section III presents the formal model divided into two contextual environments. The first considers a setting where political rewards and incentives are constant, regardless of other elites' endorsements. The second analyzes a setting where political rewards are divided among all supporters of a winning candidate. Finally, section IV includes the model results and discussion of its connection with endorsement patterns in previous elections.
